\section{Longer training and a second traffic population}
The negative short-budget comparison motivated one bounded follow-up. We fixed 20,000 updates for all four model variants and trained them again from scratch. The primary evaluation used eight separate historical dataset A validation episodes. From each episode we selected the first three plannable, strictly reconstructed starts per original class, yielding 120 requests. Each model received the same planned path and four paired noise draws. The historical dataset B validation cohort and the original 57 CNR development conditions were secondary comparisons. historical dataset A and historical dataset B share a map but differ in demand construction and eligibility.

\begin{table}[htbp]
\centering\small
\caption{Five-class macro Wasserstein discrepancy in 127.75-second travel distance (metres, lower is better). C is evaluated on the same five classes as A, D and B}
\begin{tabular}{lrrrr}\toprule
Development population & A & D & B & C\\\midrule
historical dataset B, 3,000 updates & 323.2 & 420.1 & 399.9 & 413.4\\
historical dataset B, 20,000 updates & 177.2 & 128.9 & 108.0 & 104.6\\
historical dataset A independent, 20,000 updates & 170.4 & 153.9 & 146.9 & 156.6\\
CNR, 20,000 updates & 188.0 & 145.6 & 150.9 & 151.6\\\bottomrule
\end{tabular}
\end{table}

On historical dataset B the augmented model reduced travel discrepancy by 69.2 metres relative to CNR-only training (95\% paired episode interval 14.4 to 109.8 metres). On the independent historical dataset A cohort the corresponding reduction was 23.5 metres (interval $-6.4$ to $70.7$). The two seeds gave changes of opposite sign on historical dataset A: for seed 42 the discrepancy increased by 1.8 metres, whereas for seed 43 it fell by 48.8 metres. The conditional interval captures variation over episodes; it does not account for drawing new model initializations. On the historical dataset A starts, mean time-aligned separation between two draws on the same path rose from 58.3 to 120.1 metres (five-class macro average). This measures spread within a sampled pool, without establishing the correct level of real-world variability. On historical dataset A the expanded model also lowered the discrepancy in mean absolute acceleration from 0.341 to 0.115~m/s$^2$. The apparent improvement in travel with additional SUMO examples beyond the fixed 425 was smaller, 7.1 metres on historical dataset A (interval $-16.2$ to $31.4$). These results support reporting the optimization budget alongside data scale, and leave the transfer of the travel benefit across simulation populations uncertain.

\section{The eight-class deliverable}
The trained C model takes a supported start and class and returns four sampled timed trajectories. Figure~\ref{fig:scaled-eight} shows the first supported validation request per class with a model seed fixed before evaluation. The curves display variation in motion along a common spatial opportunity within each panel. Different panels involve different roads and cannot identify the effect of switching a label on a common road.

\fig{expanded-motion.pdf}{.92}{Cumulative distance for four diffusion draws from each illustrated request after 20,000 updates. The common axes show the scale difference between walking and motorized movement. All draws for the fixed request are included}{scaled-eight}

Across both seeds, the archive contains 1,400 trajectories from 175 starts. All 1,232 motor-vehicle and bicycle trajectories passed the directed lane audit. The 168 pedestrian samples stayed within the two-metre sidewalk tolerance (0.01 m numerical allowance). The pedestrian check covers a single segment; 301 attempted starts did not meet its planning rule before the quota was reached or the search exhausted. Seven samples reached the end of their planned path, all in the bicycle class; all seven remain in the dataset. In the 3,000-update eight-class run there had been 45 caps, including 28 pedestrian cases. Geographical feasibility can therefore coexist with truncated motion. Per-class travel discrepancy also varied, ranging from 8.1 metres on the restricted pedestrian requests to 197.9 metres for buses. Bus schedules and passenger stops would be relevant conditions for a more realistic bus service.

\section{Can the inherited feedback transfer to the expanded generator?}
The calibrated G12 scorer ranked the four candidates for each B request on historical dataset A. All 960 trajectories encoded successfully. We used the requested class and the original Hilbert image channels for selection, and consulted reference futures only for evaluation.

Relative to uniform choice, the within-request speed-distribution error changed by $+0.052$ m/s (95\% paired episode interval $-0.107$ to $+0.204$); the delivered travel discrepancy changed by $-1.7$ m (interval $-60.7$ to $+46.4$). The earlier speed benefit did not replicate. Truck and delivery classes fall outside the frozen scorer’s labels.
